\section{Постановка задачи}

В данной работе рассматривается решение задачи линейного программирования с использованием симплекс-метода, а также исследование влияния абсолютной погрешности в задании правых частей ограничений на решение как прямой, так и двойственной задачи.

Целью исследования является не только нахождение оптимального решения, но и оценка его устойчивости к малым возмущениям входных данных. Для этого решается специально сконструированная задача, а затем проводится анализ чувствительности полученных решений.

\subsection{Формулировка задачи}
Рассмотрим задачу линейного программирования следующего вида:

$$
\text{Найти } x_1, x_2, x_3, x_4, x_5 \text{ такие, что }
$$

$$
Z = c_1 x_1 + c_2 x_2 + c_3 x_3 + c_4 x_4 + c_5 x_5 \to \max
$$

при ограничениях:

$$
a_{11} x_1 + a_{12} x_2 + a_{13} x_3 + a_{14} x_4 + a_{15} x_5 = b_1,
$$
$$
a_{21} x_1 + a_{22} x_2 + a_{23} x_3 + a_{24} x_4 + a_{25} x_5 = b_2,
$$
$$
a_{31} x_1 + a_{32} x_2 + a_{33} x_3 + a_{34} x_4 + a_{35} x_5 \geq b_3,
$$
$$
a_{41} x_1 + a_{42} x_2 + a_{43} x_3 + a_{44} x_4 + a_{45} x_5 \leq b_4.
$$

Условия на знак переменных:

$$x_i \geq 0, \quad i = 1, 2, 3$$

Среди коэффициентов целевой функции допускается использование нулевого коэффициента ровно один раз.

\subsection{Метод решения}

Решение задачи осуществляется с использованием симплекс-метода. % Начальное базисное решение выбирается методом искусственного базиса.

Далее проводится исследование влияния абсолютной погрешности в правых частях ограничений на полученное решение. Теоретические оценки этого влияния приводятся аналитически и проверяются численными экспериментами. Данное исследование выполняется как для прямой задачи, так и для её двойственной формы.

Практическая часть лабораторной работы реализована в виде интерактивного приложения на языке C\#, что позволяет не только получить итоговое решение, но и проследить весь процесс преобразования задачи и ход вычислительного эксперимента через графический пользовательский интерфейс.

% \section{Теоретическая часть}

% % \subsection{Метод искусственного базиса}

% % Метод искусственного базиса применяется для нахождения начального допустимого базисного решения в задачах линейного программирования, где исходные ограничения не позволяют явно выделить базисные переменные. Алгоритм включает следующие этапы:

% \begin{enumerate}
%     \item \textbf{Приведение к канонической форме:} Все неравенства преобразуются в равенства:
%     \begin{itemize}
%         \item Для ограничений типа $\leq$ добавляются переменные ($+s_i$).
%         \item Для ограничений типа $\geq$ вводятся переменные ($-s_i$).
%     \end{itemize}

%     % \item \textbf{Введение искусственных переменных:} В каждое уравнение, не содержащее базисную переменную (например, исходные равенства или преобразованные неравенства с избыточными переменными), добавляется \textit{искусственная переменная} $R_i \geq 0$.

%     % \item \textbf{Формулировка вспомогательной задачи:} Строится новая целевая функция:
%     % $$
%     % F = R_1 + R_2 + \dots + R_m \to \max
%     % $$
%     % Решение этой задачи симплекс-методом позволяет найти НДБР исходной задачи.

%     % \item \textbf{Исключение искусственных переменных:} Если в оптимальном решении $W = 0$, все $R_i$ исключаются, и полученный базис используется для исходной задачи. Если $W > 0$, исходная задача не имеет допустимых решений.
% \end{enumerate}



% \subsection{Симплекс-метод}

% Симплекс-метод — итеративный алгоритм решения задач линейного программирования, движущийся по вершинам многогранника допустимых решений. Основные компоненты:

% \begin{itemize}
%     \item \textbf{Базисные и свободные переменные:} Из $n$ переменных выбирается $m$ базисных (соответствуют ненулевым значениям), остальные — свободные (равны нулю).
    
%     \item \textbf{Симплекс-таблица:} Матрица коэффициентов системы ограничений и целевой функции:
%     \begin{itemize}
%         \item Строки соответствуют базисным переменным.
%         \item Столбцы — всем переменным и правым частям.
%         \item Последняя строка содержит коэффициенты целевой функции (индексная строка).
%     \end{itemize}

%     \item \textbf{Критерии выбора:}
%     \begin{itemize}
%         \item \textit{Ведущий столбец:} Для максимизации выбирается столбец с наименьшим отрицательным коэффициентом в индексной строке.
%         \item \textit{Ведущая строка:} Определяется по минимальному отношению правой части к положительному элементу ведущего столбца (правило минимального отношения).
%     \end{itemize}

%     \item \textbf{Правило прямоугольника:} Пересчет таблицы выполняется по формулам:
%     $$
%     \Delta_{ij} = \frac{a_{ik} \cdot a_{lj}}{a_{lk}}
%     $$
%     $$
%     a'_{ij} = a_{ij} - \Delta_{ij}
%     $$
%     где $a_{lk}$ — ведущий элемент.
% \end{itemize}

% Алгоритм завершается, когда в индексной строке отсутствуют отрицательные коэффициенты — достигнута оптимальная вершина.
% \newpage
% \subsection{Блок-схема симплекс-метода}

% % \begin{enumerate}
% %     \item Начало.
% %     \item Приведение задачи к канонической форме.
% %     \item Построение начальной симплекс-таблицы.
% %     \item Проверка индексной строки на оптимальность:
% %     \begin{itemize}
% %         \item Да → переход к шагу 8.
% %         \item Нет → выбор ведущего столбца.
% %     \end{itemize}
% %     \item Выбор ведущей строки по правилу минимального отношения.
% %     \item Пересчет таблицы 
% %     \item Повтор шагов 4-6 до выполнения условия оптимальности.
% %     \item Вывод результатов: значения переменных, оптимальное значение целевой функции.
% %     \item Конец.
% % \end{enumerate}

% % \begin{itemize}
% %     \item \textit{Особенность для метода искусственного базиса:} Перед шагом 2 выполняется введение искусственных переменных и решение вспомогательной задачи.
% %     \item \textit{Вырожденность:} Если минимальное отношение достигается для нескольких строк, возникает вырожденность, требующая специальных стратегий (например, правило Бленда).
% % \end{itemize}

% \begin{figure}[htbp]
%     \centering
%     \includegraphics[width=1.0\textwidth]{blok_shema.png}
%  \end{figure}
